\scanpage{36}
\mspage{35}

Soit $\deleted{\unsure{\mathcal F_{\mathrm{max.}}}\ \unsure{\mathcal F^{*}}}$ $\Phi$ l'ens des parties de $M$ qui sont a) fermées pour $\leq$ et b) formées d'éléments deux à deux compatibles. (au sens de (1) p.~29). Munissons $\Phi$ de la relation $\leq\ \overset{\mathrm{def}}{=}$ l'inclusion, et $\ll$ définie par $F\ll G\ \overset{\mathrm{def}}{\Longleftrightarrow}\ \forall X\in F,\exists\,Y\in G$ avec $X\ll Y$. Alors $\Phi$ est un atlas \deleted{admettant} satisfaisant $\mathrm{\deleted{\unsure{M|A}}At\,L\,1}$, $\mathrm{At\,L\,2}$, $\mathrm{At\,L\,3}$ et $\mathrm{At\,CL(\unsure{\mathrm{spac}|\mathrm{spa}})}$, et admettant les $\widetilde X=M_{\leq\unsure{X|x}}$ $\deleted{\unsure{\{\{X\}\}|\{\{x\}\}}}$ $(X\in M)$ comme \deleted{objets} figures élémentaires, \deleted{les $\{\{x\}\}$ \unsure{pour}} et \deleted{\unsure{qui}} induisant sur $M$ les relations $\leq$, $\ll$ de départ, et sur $\unsure{L|\mathcal L}$ \inserted{(pour la \deleted{\unsure{dis}} disjonction des figures)} la relation $\parallel$.

\deleted{Tout \unsure{autre} \unsure{les atlas|tel atlas}} \inserted{Tout atlas correspondant à ces} $\unsure{\Rightarrow|\mathord{=}}$ données s'identifie à \unsure{un|son} \uline{sous-atlas} \uline{spacieux} de $\Phi$\unsure{, »|,} \inserted{\unsure{i.e.|i.e}} \deleted{(parties $\unsure{\mathcal F|F}$ de $\Phi$} à une partie
\[
 \mathcal F\subset\Phi
\]
satisfaisant

(\deleted{\unsure{a|+}}i) $\forall\,F\in\mathcal F$, et $G\leq F$ \deleted{\unsure{sous-figure}} \inserted{(partie fermée} de $F,\leq$) \unsure{on a} $G\in\mathcal F$ et

(ii) $\forall\,F,G\in\mathcal F$, Si $F\mathrel{\unsure{\ooalign{\hfil$>$\hfil\cr\hfil$<$\hfil\cr}|\gtrless}}G$ i.e. $F\cup G\in\Phi$ (i.e. \deleted{\unsure{tous}}\unsure{tout|tous} éléments de $F_{\unsure{a|x}}$ est comp. à tout \unsure{élément|élmt|élémts} de $G$) $F\cup G\in\mathcal F$.
